\chapter{Remaining FYP Plan}
\label{ch:plan}
% ============================================================

\section{Remaining Scope}
Objectives O4--O6: detection model, federated training and comparative evaluation,
and attributed alerting driven by real model output.

\section{Remaining Modules and Features}
M4 Federated Detection Model; M5 Zone Participation and Aggregation; M6 Alerting and
Presentation; the redundant-topology and IEEE 802.1CB FRER extension to M2 (topology
and attack injection) and M3 (feature extraction); plus corrective work on M2 and M3
(DEF-01, DEF-02).

\section{Remaining Iterations}
\label{sec:remaining-iterations}

\textbf{Iteration 03: Detection and decision gate. Complete.}
The cross-zone sibling residual was added, a centralized detector trained, and the
feature ablation completed. The decision gate is passed: cross-zone features exceed
within-zone features by $+0.43$ to $+0.49$ macro-F1 across three models, so the
federated work proceeds as planned.

\textbf{Iteration 03b: Dataset depth.}
Extract the \texttt{SyncStarvation} configuration to give a genuine multiclass problem.
Sweep the injected attack intensity downwards to establish the detection boundary
against the benign envelope. Introduce differentiated per-zone traffic profiles
(DEF-02).

\textbf{Iteration 03c: Redundant topology and FRER.}
The topology half of this iteration is complete in design: the five-zone ring
(Section~\ref{sec:ref-topology}) is generated and passes its structural checks, including the
single-link-cut partition test that establishes disjoint paths. What remains is to run
the benign baseline on it and record the four behavioural checks listed in
Section~\ref{sec:topology-validation}, then configure IEEE 802.1CB sequence generation at the talkers and
sequence recovery at the listeners for the redundant streams. Implement
stream-redundancy manipulation as a third attack class: selective suppression or
delay of one member of a replicated stream in a single target zone. Add the FRER
feature family to \texttt{extract\_features.py}, namely inter-member arrival skew,
sequence-number deviation between redundant members, and per-member inter-arrival
cadence, together with their cross-zone sibling residuals, computed by the same
window-outer, zone-inner construction already used for the gPTP features. Exit
criterion: the FRER residual features reach macro-F1 $> 0.90$ on the redundancy
attack class while the corresponding within-zone features remain at chance, which
would show the residual principle to be a property of the zonal architecture rather
than of the synchronisation plane alone. This iteration is independent of
Iteration~03b and may run in parallel, since it touches the topology and the feature
extractor rather than the detector. Migrating the gPTP attack configurations from the
four-zone tree to the five-zone ring is a prerequisite for both, and is the point at
which the Chapter~\ref{ch:testing} results are re-established on the reference topology.

\textbf{Iteration 04: Federation and evaluation.}
Federated client/server implementation; local-only, centralized-pooled and FedAvg
arms on identical splits; zone attribution evaluation; detection latency; alert
surface driven by model output; rule-based fallback; monitoring-interference
measurement.

\section{Implementation Plan}
Iteration 03 proceeds in strict order: labelling correction, heterogeneity, third
attack class, control arm, comparison. The ordering is deliberate, because each step is a
precondition for interpreting the next. Iteration 03c runs in parallel on the
topology and extraction path: redundant links, 802.1CB sequence generation and
recovery, redundancy attack injection, then the FRER feature family and its
residuals. Iteration 04 proceeds in order: local-only
baseline, centralized upper bound, FedAvg, then the remaining evaluation dimensions.

\section{Testing and Validation Plan}
\label{sec:validation-plan}

\begin{longtable}{p{0.28\textwidth} p{0.66\textwidth}}
\toprule
\textbf{Evaluation dimension} & \textbf{Metrics} \\
\midrule
\endhead
Detection quality & Per-class precision, recall, F1; macro-F1; confusion matrix;
detection rate at fixed false-positive rate \\
Feature-set comparison & Macro-F1 of protocol-native versus generic traffic features,
identical model and splits, reported across seeds \\
Operational detection & Detection latency in synchronisation intervals; false alarms
per vehicle-hour \\
Attribution & Zone attribution top-1 accuracy, reported separately from detection \\
Federated behaviour & Worst-client macro-F1; centralized-to-FedAvg gap; convergence
rounds; update volume per round \\
Baselines & Fixed threshold on raw offset; CUSUM on sibling residual; LOF (for
comparability with published gPTP work) \\
TSN behaviour & Protected-stream mean and 99th-percentile latency, monitoring enabled
versus disabled (NFR-08) \\
Reproducibility & Results across fixed seeds, held-out attack episodes and held-out
configurations \\
\bottomrule
\end{longtable}

Accuracy is excluded as a headline measure throughout, owing to the class imbalance
reported in Section~\ref{sec:test-results}. Train/test separation is performed at run and seed
granularity (integrity rule R3).

\section{Deployment Plan}
\label{sec:deployment-plan}
The deliverable is a reproducible research and validation environment: pinned
dependencies, documented commands, sample configurations, seeds, and exported
results, executable by a non-author on a clean environment.

Two consumers of detector output are defined. Within the vehicle, the integrity
verdict (class, affected branch, confidence) is the decision input that
fault-tolerant timing mechanisms require in order to fail over or enter holdover; the
relevant standardised machinery provides the failover mechanism but leaves the
domain-quality decision unspecified. Off-vehicle, attributed alerts populate the vSOC
incident queue, where per-vehicle anomalies are triaged against fleet-wide
correlation.

\section{Tentative Timeline}
The project proposal was defended on 11 September 2026 and this mid evaluation falls
in the week of 13 October 2026. The final FYP-1 evaluation is scheduled for December
2026, which is the deadline the sequence below is sized against. The sequence is
ordered by dependency rather than by date. Iterations 03b and
03c are independent of each other and may run in parallel; Iteration 04 depends on
both, because the federated comparison is meaningless until the clients differ
(03b) and incomplete until the second attack family exists (03c).

\begin{longtable}{p{0.16\textwidth} p{0.40\textwidth} p{0.34\textwidth}}
\caption{Tentative FYP-2 timeline}\label{tab:timeline}\\
\toprule
\textbf{Iteration} & \textbf{Work} & \textbf{Exit condition} \\
\midrule
\endfirsthead
\toprule
\textbf{Iteration} & \textbf{Work} & \textbf{Exit condition} \\
\midrule
\endhead
03b &
Extract \texttt{SyncStarvation}; sweep attack intensity downward; introduce
differentiated per-zone traffic profiles &
A genuine multiclass dataset exists, and the per-zone divergence of DEF-02 is
measured rather than assumed \\
\addlinespace
03c &
Benign baseline run on the ring topology; 802.1CB sequence generation and recovery;
redundancy attack injection; FRER feature family &
The four behavioural criteria of Section~\ref{sec:topology-validation} are met, and FRER residuals reach
macro-F1 $> 0.90$ while the matching within-zone features remain at chance \\
\addlinespace
04 &
Local-only, centralized-pooled and FedAvg arms on identical splits; zone
attribution; detection latency; alert surface &
SC-4 is answered either way: the federated arm is measured against both baselines
on the same splits, and a null result is reported as a result \\
\addlinespace
Final &
Conformal calibration; monitoring-interference measurement; write-up &
NFR-08 measured rather than carried over; final report and defence \\
\bottomrule
\end{longtable}

Iterations 03b and 03c fall between this mid evaluation and the final FYP-1
evaluation in December 2026. Iteration~04 and the final items belong to FYP-2.

% [[FILL]] If your department publishes dated FYP-2 milestones, add a date column to
% the table above. The ordering and the exit conditions do not change; only the dates.

\section{Expected Final Deliverables}
Reproducible zonal TSN simulation with attack injection; labelled protocol-native
feature dataset; trained detection model with federated and centralized training
paths; comparative evaluation report covering detection, attribution, latency and
federated behaviour; operator-facing alert surface; documented negative results and
scope boundaries; complete source, configuration and handover package.

% ============================================================